\subsection{Resultados e análise}

As saídas híbridas coincidiram com a saída sequencial nas 144
execuções paralelas, incluindo os aquecimentos. Em uma validação separada, os 50 casos sintéticos também coincidiram
com o oráculo e com o sequencial. As 15 entradas inválidas foram
rejeitadas. Os checksums das cargas oficiais foram
\begin{itemize}
  \item pequena: \texttt{17073017104646724864}
  \item média: \texttt{17187382406529706254}
  \item grande: \texttt{14825089465781038537}
\end{itemize}

A configuração $2\times4$ em dois nós apresentou a menor média nas
três cargas (Tabelas~\ref{tab:lasdpc-pequena} a~\ref{tab:lasdpc-grande}).
Na carga grande, o tempo médio passou de 13,10~s no sequencial para
2,05~s, com speedup de 6,39 e eficiência de 79,8\%. Nessa
configuração, a eficiência na carga pequena foi de 62,3\%.

Com dois ranks e duas threads por rank, calculamos em cada rodada a
diferença de tempo entre dois nós e um nó. Na carga pequena, a
diferença média foi de $0{,}0173$~s, com IC de 95\% de
$[0{,}0109; 0{,}0236]$~s. Na média e na grande, os tempos em dois
nós foram menores. As diferenças médias foram $-0{,}0210$~s
e $-0{,}0907$~s, respectivamente. Essa comparação mede o efeito
da colocação como um todo. Alterar a colocação muda tanto o caminho
de comunicação quanto os recursos locais que os ranks compartilham.

A Figura~\ref{fig:tb2-desempenho} reúne tempo, speedup e eficiência.
Nos rótulos, \texttt{one} identifica o \texttt{hal01} e
\texttt{two} identifica \texttt{hal01} e \texttt{hal02}.
